%=====================================================================
\chapter{Answers to the Checkpoints}\label{app:answers}
%=====================================================================
\normalfigures

Three answers per chapter, in order. Attempt the checkpoint from memory before
reading these, and argue with them where you disagree --- several are judgement
calls rather than facts, and they say so.

% The slide pipeline parses this appendix: \ansch{<number> --- <title>} must be
% followed immediately by an ans environment of exactly three items.
\newcommand{\ansch}[1]{\subsection*{\color{primary}Chapter #1}}
\newenvironment{ans}
  {\begin{enumerate}[leftmargin=*, itemsep=4pt, topsep=3pt,
                     label=\textbf{\arabic*.}]}
  {\end{enumerate}}

%---------------------------------------------------------------------
\ansch{1 --- What E-Commerce Is, and What It Changed}
\begin{ans}
  \item \textbf{Electronic business, not electronic commerce.} Viewing results
        and downloading transcripts are business processes conducted over a
        network, but no value changes hands through the system --- the
        transaction happens at a bank counter. The moment fee payment moves
        online it becomes e-commerce, and the university acquires every problem
        to the right of the dashed line in \dsref{fourfunctions}: a payment
        that may have succeeded while the network dropped, a receipt that must
        exist exactly once, a refund path.
  \item Many are acceptable. One good answer: \textbf{refusing to sell the last
        item when it is visibly damaged.} In software this requires a stock
        record that distinguishes ``held'' from ``sellable'', somebody whose
        job is to update it, a rule for what the storefront does when the two
        disagree, and a decision about what the customer sees --- out of stock,
        or a quiet substitution? The shopkeeper does all of that in half a
        second without recording that a decision was made, which is exactly why
        it is forgotten when the shop goes online.
  \item \textbf{Two uncounted costs:} the acquisition cost of finding customers
        who currently arrive free (\chref{models} puts this at Rs.~1{,}200 an
        order for a typical small shop, which can exceed the commission), and
        the labour of running the four functions themselves --- payment
        integration, delivery arrangements, customer service, returns. The
        commission was buying those, not just the listing. \textbf{What they
        gain that is not money:} the customer relationship. They learn who
        bought, can contact them again, and cannot be delisted. For a business
        with repeat custom that is frequently worth more than the 15\%.
\end{ans}

%---------------------------------------------------------------------
\ansch{2 --- Business Models and Revenue Models}
\begin{ans}
  \item \textbf{Business model:} a marketplace connecting restaurants with
        diners, with fulfilment performed in-house --- three and a half of the
        four functions. \textbf{Revenue models: three at once} ---
        \emph{commission} from restaurants, a \emph{service or delivery fee}
        from customers, and \emph{advertising} in the form of promoted
        placement. \textbf{Most profitable per unit: the advertising.} It has
        almost no marginal cost, where commission is earned against riders and
        support that must be paid for, and the delivery fee rarely covers the
        rider. Promoted placement is close to pure margin, which is why these
        platforms push it so hard.
  \item Viable when $\text{LTV} > \text{CAC}$, that is when
        $400 \times n > 1{,}500$, so when the average customer orders
        $n > 3.75$ --- \textbf{at least four times} --- before they stop.
        \emph{What to measure:} the repeat rate and the time over which it
        happens. Four orders over six months is a business; four orders over
        six years is not, because the cash was spent today and returns over a
        period during which the venture has to survive.
  \item B2C platforms store \textbf{one price per product} --- a column on the
        product row. B2B pricing is a \emph{relation} between customer and
        product, often with quantity breaks and expiry dates. That is not a
        setting you have failed to find; it is a different cardinality, and
        making it work means either a different platform or an extension that
        intercepts every price lookup in the system (\chref{modules} shows how
        much work that is). The assumption that breaks is ``a product has
        \emph{a} price''.
\end{ans}

%---------------------------------------------------------------------
\ansch{3 --- Auctions, Portals and Communities}
\begin{ans}
  \item \textbf{Bid Rs.~15{,}000} --- your true valuation. You win, and you pay
        the second-highest bid, \textbf{Rs.~13{,}500}, gaining Rs.~1{,}500.
        Bidding Rs.~18{,}000 does not improve this outcome at all: you still
        pay Rs.~13{,}500, because the price is set by the others. What it does
        is expose you to winning when somebody bids Rs.~16{,}500, at which
        point you pay Rs.~16{,}500 for a thing worth Rs.~15{,}000 to you. The
        overbid can only ever buy you an auction you did not want.
  \item It \textbf{never closes}. The auction stays open indefinitely; the
        winning bidder is never notified, the seller is never paid, and the
        listing accumulates further bids whenever somebody wanders past. The
        seller sees an item that appears to be permanently ``ending soon'' and
        a buyer who has given up. Worse, the failure is \emph{silent} and
        inversely proportional to interest: popular auctions close correctly
        and unpopular ones do not, so the bug is invisible in testing, where
        you load every page you create.
  \item Any two of: \textbf{be a seller yourself} until real ones arrive
        (cost: you take on stock and inventory risk, and you compete with the
        sellers you are recruiting); \textbf{narrow the market} to one city or
        one category (cost: you cap your growth and may pick the wrong niche);
        \textbf{give sellers a tool worth having without buyers} (cost: you are
        now building and supporting a second product); \textbf{list supply
        gathered elsewhere} (cost: legally risky and reputationally
        expensive if discovered).
\end{ans}

%---------------------------------------------------------------------
\ansch{4 --- Planning an E-Commerce Framework}
\begin{ans}
  \item (a)~\textbf{Functional}, and it \emph{does not} discriminate --- all
        four options reset passwords. (b)~\textbf{Functional}, and it
        \textbf{discriminates sharply}: per-customer pricing eliminates most
        hosted platforms and the marketplace route outright, and is the single
        most common reason to choose self-hosted or bespoke. (c)~\textbf{
        Non-functional}, and it discriminates \emph{in the opposite direction}
        to what students expect --- surviving a twentyfold spike is something a
        large hosted platform does effortlessly and your single virtual server
        does not.
  \item \textbf{Hosting} (Rs.~324{,}000 in the worked example),
        \textbf{setup and theming labour} (Rs.~175{,}000 plus
        Rs.~100{,}000 of custom modules), and above all \textbf{maintenance
        labour} (Rs.~360{,}000). Their stated saving of Rs.~792{,}000 is
        therefore closer to Rs.~339{,}000 --- still a saving, but less than
        half what they claimed, and it vanishes entirely if maintenance takes
        eight hours a month instead of four.
  \item Many are acceptable. A clear one: \textbf{a home baker selling six
        kinds of cake to one city, taking orders through Instagram.} Self-hosting
        would mean paying for a server and maintaining it to serve perhaps
        thirty orders a week, with nobody to patch it --- which is not a
        cheaper shop but an unpatched one. \textbf{Choose marketplace-only or a
        cheap hosted platform}: the business's scarce resource is the baker's
        time, the volume is low, nothing about the pricing is unusual, and
        every hour spent on a server is an hour not spent baking.
\end{ans}

%---------------------------------------------------------------------
\ansch{5 --- The Stack, and Your Shop in One Command}
\begin{ans}
  \item \textbf{Working:} the operating system and the web server --- a 500 is
        generated and served by Apache, so the request reached it and a
        response came back. \textbf{Failed:} the application layer; PHP raised
        an error and died before producing a page. The database is
        \emph{unknown} --- it may be the cause (the application could not
        connect) or entirely healthy. \textbf{Look first} at the PHP error log,
        with \texttt{docker compose logs shop}; the message is never on the
        screen, because \texttt{display\_errors} is off, and \chref{security}
        is about why you want it that way.
  \item \texttt{depends\_on} orders container \emph{starts}, and a container
        that has started is not a database that will accept a login --- MariaDB
        opens its TCP port several seconds before it finishes initialising. The
        shop therefore starts, runs the installer, and the installer is refused
        by a database that is listening but not ready. \textbf{The failure is
        intermittent and worse on a fast machine}, because the faster the shop
        container starts, the more certainly it arrives early. The
        \texttt{service\_healthy} condition waits for the database's own
        healthcheck to pass, which tests a login rather than a port.
  \item \textbf{Lost:} both named volumes --- the entire database (products,
        categories, orders, customers, settings) and the shop's installed
        files, including any module you had written in place. \textbf{Not
        lost:} anything in the repository --- the Compose file, the Dockerfile,
        \texttt{code/}, and any module kept there rather than edited inside the
        container. \textbf{What would have made it minor:} a database dump
        taken on a schedule and kept outside the volume, plus the discipline of
        writing modules in \texttt{code/modules/} and mounting them, so that
        the only thing destroyed is state you can regenerate. \chref{deploy}
        makes both of these routine.
\end{ans}

%---------------------------------------------------------------------
\ansch{6 --- Products, Categories and Their Tables}
\begin{ans}
  \item \textbf{Move to \texttt{oc\_product\_description}:} the product name
        and the meta description. Both are \emph{words}, and words differ by
        language. \textbf{Stay in \texttt{oc\_product}:} price, weight and
        stock quantity --- a kilogram is a kilogram in Urdu, and a shop that
        let the price depend on the reader's language would eventually charge
        two different prices for one product. The test is the first of the four
        questions: \emph{does this fact depend on who is reading?}
  \item \textbf{Buys:} a single-query answer to ``every category beneath this
        one'', so a category page costs one query instead of a loop whose
        length depends on how deep the tree happens to be.
        \textbf{Costs:} the same fact is now stored twice, so the two can
        disagree, and every write that changes the shape of the tree must
        rewrite the derived rows. \textbf{The operation that makes it visible:}
        \emph{moving a category}. Changing one \texttt{parent\_id} invalidates
        the path rows of that category and of everything beneath it, so the
        admin panel must rewrite them all --- which is why moving a category
        with many descendants is noticeably slower than renaming one.
  \item Two of several good answers. \textbf{(a)~The database was never the
        bottleneck}: if the database is 20\% of the page, halving it improves
        the page by 10\%, which can be smaller than the run-to-run variance.
        \textbf{(b)~The expensive thing is the \emph{number} of queries, not
        the cost of any one} --- the measurement in this chapter found 146
        queries for one page, and an index helps exactly one of them.
        \smallskip
        \textbf{To tell them apart:} turn on the query log with
        \texttt{long\_query\_time=0}, request the page once, and read
        \texttt{COUNT(*)} and \texttt{SUM(query\_time)} from
        \texttt{mysql.slow\_log}. A large total with few queries is (a); a
        small total spread over a hundred queries is (b). That single
        measurement separates them, and it is the one in
        \texttt{m06\_catalogue.py}.
\end{ans}

%---------------------------------------------------------------------
\ansch{7 --- Product Variations, Options and Uploads}
\begin{ans}
  \item \textbf{The photograph is an option} --- specifically a \texttt{file}
        option. It is supplied by the customer, it is different every time, and
        the shop holds no stock of it. \textbf{The size is a variant}, because
        a small mug and a large mug are two different objects on two different
        shelves, each with its own cost, its own barcode and its own count.
        \textbf{What the shop holds in stock:} blank mugs in two sizes. Nothing
        else --- which is exactly why the photograph cannot be a variant: there
        is no such thing as stock of ``the mug with Ahmed's photograph on it''
        until Ahmed orders one.
  \item \textbf{Option value rows: $4 + 3 + 2 = \mathbf{9}$}, plus three rows
        in \texttt{oc\_product\_option}. \textbf{As variants: $4 \times 3 \times
        2 = 24$ combinations, so \textbf{25} \texttt{oc\_product} rows} ---
        twenty-four plus the master --- and 25 description rows, 25 store rows
        and 25 category rows beside them. Nine against a hundred, for the same
        catalogue.
  \item \textbf{Passes:} the filename sanitiser (nothing illegal in
        ``logo.svg''), the length check, the extension allow-list
        (\texttt{svg} is on the stock list), and the MIME allow-list --- it
        claimed \texttt{image/png}, which is allowed, and nothing verifies that
        claim. \textbf{Fails:} nothing. All five checks pass.
        \smallskip
        \textbf{For it to harm a visitor}, the file must be \emph{served to
        them from your origin and rendered}. Two things would have to be true:
        the attacker must learn the stored URL, which carries a 32-character
        random token and so is not guessable --- but it \emph{is} shown to
        whoever uploaded it and may be forwarded to an administrator who opens
        it --- and the server must serve it with
        \texttt{Content-Type: image/svg+xml} rather than as a download, which
        \chref{stack}'s stock install does. The realistic victim is therefore
        the administrator reviewing the order, and the realistic fix is to drop
        \texttt{svg} from the list and serve uploads with
        \texttt{Content-Disposition: attachment} (\chref{security}).
\end{ans}

%---------------------------------------------------------------------
\ansch{8 --- Themes, Layout and the Customer's Experience}
\begin{ans}
  \item \textbf{Best case for compressing the photographs:} the images are
        $713 - 597 - 54 = $ about \textbf{61\,KB}. Halving them saves 30\,KB,
        which is \textbf{4\%} of the page --- and that is the \emph{best} case,
        because you cannot compress them to nothing. \textbf{Propose instead:}
        minify \texttt{bootstrap.css}, which is 265\,KB of the 597 and has no
        minified copy in the tree. That single change saves roughly 240\,KB,
        about \textbf{34\%} of the page --- eight times the whole image budget
        --- and changes nothing a customer can see. Then remove the unused
        parts of Bootstrap and FontAwesome, and only then look at the images.
  \item \textbf{Layout change} if a ``trust badge'' module already exists:
        \adminpath{Design, Layouts, Product}, add the module to
        \uiel{Content Bottom}, set its sort order. No code, visible to the next
        person, survives upgrades. \textbf{Template override} if it must sit
        \emph{immediately under the button}, inside the markup the theme
        generates: there is no position there, so you override
        \texttt{product/product.twig} in your child theme and add it. The
        general rule: a position that exists is a layout change; markup that
        must differ is a template change.
  \item \textbf{Two reasons it is lower.} (a)~\textbf{It follows only assets
        named in the HTML.} Fonts are requested by the CSS once the browser
        parses it, so none of the FontAwesome \texttt{.woff2} files --- 248\,KB
        in that directory --- are counted. (b)~\textbf{It does not execute
        JavaScript}, so anything a script loads afterwards (a tracking pixel,
        a lazy-loaded image, a carousel's later slides) is invisible to it.
        \smallskip
        \textbf{Which it could fix:} the first, by parsing the CSS for
        \texttt{url()} references and following them --- a dozen lines. The
        second it cannot fix, because fixing it means running a browser, which
        is a different kind of tool. Knowing which limitation is laziness and
        which is fundamental is the useful judgement here.
\end{ans}

%---------------------------------------------------------------------
\ansch{9 --- The Shopping Basket, and Where State Lives}
\begin{ans}
  \item \textbf{What would change:} the two queries per request that the
        \texttt{db} engine spends reading and writing the session record. On
        this host that was worth about 20\% of the page --- 89\,ms to 71\,ms
        --- when those two queries were removed altogether by the
        \texttt{file} engine, and Redis would remove them from the database in
        the same way. \textbf{What would not change: the basket.} It is rows in
        \texttt{oc\_cart}, and the \texttt{Cart rows} column was 1 under both
        engines. \textbf{The number that settles it:} the session record is
        \textbf{18 bytes} and holds \texttt{\{"currency":"USD"\}}. Whatever you
        move it to, you are moving eighteen bytes. The right place to look for
        basket performance is the six cart queries an empty basket costs, not
        the store the session sits in.
  \item \textbf{Where they came from:} a session row is written on every
        request, before the visitor has asked for anything, so every page view
        by every crawler, uptime check and measurement harness created one. The
        645 rows here are the by-product of writing this handout.
        \textbf{Why they are not removed on departure:} nothing tells a server
        that a visitor has left --- HTTP has no ``goodbye''. Rows can therefore
        only expire on a clock, and \texttt{config\_session\_expire} is 86{,}400
        seconds. \textbf{What to measure first:} the proportion of sessions
        that ever acquire a basket or a login. If almost none do, the expiry is
        not the problem and robots are; shortening the expiry would then also
        shorten the life of the few baskets that matter, which is the opposite
        of what was wanted. Measure before tuning the number that is easiest
        to reach.
  \item \textbf{What each removes:} \texttt{down} removes the containers ---
        the running processes and their writable layers. The named volumes
        \texttt{dbdata} and \texttt{shopdata} are separate objects and are left
        alone, so the database files survive and the basket with them;
        \texttt{up -d} attaches the new containers to the same volumes and the
        shop resumes mid-visit. \texttt{down -v} additionally deletes those
        volumes, which is where every byte of durable state lives, so the
        entrypoint finds an empty database and reinstalls from scratch.
        \textbf{One other thing it destroys:} every product, category, order,
        customer and setting you have made since \chref{stack} --- in this
        measurement the shop came back in 17 seconds looking perfectly
        healthy, with all of it gone. Uploaded images go with the
        \texttt{shopdata} volume at the same time.
\end{ans}
